\chapter{Arhitectura Aplicației}
\label{cap:arhitectura}

\section{Arhitectura Generală Client-Server}

Aplicația SMS urmează o arhitectură clasică client-server:

\begin{itemize}
  \item \textbf{Frontend (Client):} Angular SPA, rulează în browser
  \item \textbf{Backend API:} PHP REST API, pe server web
  \item \textbf{Database:} MySQL, accesibilă doar din backend
  \item \textbf{Communication:} HTTP/HTTPS cu JSON
\end{itemize}

\begin{figure}[H]
  \centering
  % \includegraphics[width=0.85\textwidth]{1_ARHITECTURA_SISTEM.drawio.png}
  \caption{Arhitectura Sistem – Client-Server}
  \label{fig:arch}
\end{figure}

\section{Structura Frontend-ului (Core, Shared, Features)}

Frontend-ul este organizat pe module:

\begin{itemize}
  \item \textbf{AppModule:} root module, routing principal
  \item \textbf{CoreModule:} servicii singleton (AuthService, HttpService)
  \item \textbf{SharedModule:} componente refolosibile (Navbar, Sidebar)
  \item \textbf{Feature Modules (lazy-loaded):} StudentModule, CourseModule, AttendanceModule, ExamModule
\end{itemize}

\section{Organizarea Modulelor și Componentelor}

\begin{figure}[H]
  \centering
  % \includegraphics[width=0.8\textwidth]{3_FRONTEND_COMPONENTS.drawio.png}
  \caption{Component Tree – Ierarhia Componentelor}
  \label{fig:components}
\end{figure}

Fiecare modul conține: Components, Services, Models (TypeScript interfaces), Routing.

Pattern: \textbf{Smart Components} apelează servicii, \textbf{Presentational Components} doar afișează.

\section{Fluxul Datelor (Component → Service → API → UI)}

\begin{enumerate}
  \item User interacționează cu Component
  \item Component apelează Service method
  \item Service apelează HttpClient
  \item HttpClient trimite request la API (cu JWT header din JwtInterceptor)
  \item API procesează, validează, salvează în DB
  \item API returnează JSON
  \item Service emite datele via Observable
  \item Component subscribe și actualizează UI
\end{enumerate}

\section{Modelul de Date (ERD)}

%\begin{figure}[H]
%  \centering
%  \includegraphics[width=0.95\textwidth]{erd-placeholder.png}
%  \caption{Entity-Relationship Diagram}
%  \label{fig:erd}
%\end{figure}

Relații principale: 1:N și M:N între tabele (teachers-courses, students-groups, sessions-attendance, etc.).

\section{Diagrame Relevante}

\begin{figure}[H]
  \centering
  % \includegraphics[width=0.85\textwidth]{2_USE_CASES.drawio.png}
  \caption{Diagrama Cazurilor de Utilizare}
  \label{fig:usecases}
\end{figure}